\section{The Lexicon}

Every theory needs a small set of words that mean exactly one thing. This section fixes them. The names here are authoritative for the rest of the paper, and each name is chosen so that it reads cleanly both as ordinary English and as a symbol in the math.

The whole picture is one sentence. The universe is a \emph{vibe mesh}, vibes vibing. Broken down more carefully, it is a \emph{flow}, eight atoms running as one.

Two words sit above the eight, the substance and the whole. The substance is the \textbf{vibe}, the one stuff that everything is made of. The whole is the \textbf{flow}, the eight atoms taken together as a single thing in motion, written $F$. The vibe is what the parts are made of. The flow is what the parts add up to. Between them lie the eight atoms, and nothing else is in the base.

\subsection{The eight atoms}

The base is exactly eight atoms. Each is a four-letter word, and each starts with a different letter, so each doubles as a clean single-letter symbol for the algebra. Together they answer the universal questions a law has to answer, a where, a what, a how, and a when, with the how split in two because two different things change each step, the state and the space.

\begin{table}[ht]
\centering
\begin{tabular}{@{}lllll@{}}
\toprule
atom & symbol & what it is & category \\
\midrule
\textbf{mesh} & $m$ & the whole $\{3,4,3,4\}$ tessellation & where (space) \\
\textbf{dock} & $d$ & the geometric cell, a here-now & where (space) \\
\textbf{site} & $s$ & one of the $24$ directions out of a dock & where (space) \\
\textbf{tone} & $t$ & the vibe value, ternary $\{-1,0,+1\}$ & what (substance) \\
\textbf{lean} & $l$ & the value order $-1 < 0 < +1$ & what (substance) \\
\textbf{knit} & $k$ & collide then stream, the state law & how (state) \\
\textbf{wake} & $w$ & the growing edge, new docks at the frontier & how (space) \\
\textbf{beat} & $b$ & the discrete step, time & when (time) \\
\bottomrule
\end{tabular}
\caption{The eight atoms of the base. Eight four-letter words, eight unique first letters $m\,d\,s\,t\,l\,k\,w\,b$, assembled into the whole $F = $ flow.}
\label{tab:atoms}
\end{table}

\begin{axiom}[The eight atoms]
The base of the theory is exactly the eight atoms mesh, dock, site, tone, lean, knit, wake, and beat. Nothing else is fundamental. Every other term in the paper either emerges from these eight or names an aspect of them.
\end{axiom}

\subsection{The eight, defined}

Each atom is defined once, here, and used with this meaning everywhere after.

\subsubsection{mesh}

The \textbf{mesh} is the whole discrete space, the $\{3,4,3,4\}$ hyperbolic honeycomb. It is finite at any single beat and forever growing. It is a weave of docks. When we say the universe is a vibe mesh, this is the mesh we mean. In a phrase, the mesh is a weave of docks.

\subsubsection{dock}

A \textbf{dock} is one cell in the mesh, a here-now, the geometric unit of place. It is a place of transit. Each beat the tones inside a dock interact, then move on to neighbour docks. The dock is where the action of one step happens.

\subsubsection{site}

A \textbf{site} is one of the $24$ directions out of a dock, the directed slot where a tone lives. Space nests. The mesh contains docks, and a dock contains exactly $24$ sites. There is no rest slot, no twenty-fifth site, no still home. A tone sits in one site and moves along it.

\begin{definition}[Dock and site]
A \emph{dock} contains exactly $24$ \emph{sites}, one per local direction of the $\{3,4,3,4\}$ cell, and no others. There is no rest slot. The $24$ together form one hyper-color, the full state a dock can hold at a beat.
\end{definition}

\subsubsection{tone}

A \textbf{tone} is the vibe value carried by a site, a ternary number in $\{-1, 0, +1\}$. Read experientially, the three values are pain, peace, and pleasure. The tone is the only attribute a site carries. A site \emph{has} a tone, and its tone can change while the site itself persists. This keeps two things apart. The bearer is the site. The value is the tone. So the vibe, the substance, stays distinct from the tone, the value it takes. The single quale of the theory is one tone.

\subsubsection{lean}

The \textbf{lean} is the value order, $-1 < 0 < +1$. The tones lean. That tilt does two jobs. It gives charge, the signed sum of the tones, and it gives the create move its direction, since peace leans out into a balanced pair. The lean is the atom of order among the values, pain below, pleasure above, peace at the hinge.

\subsubsection{knit}

The \textbf{knit} is the law of how the state changes each beat. It has two phases, collide then stream. In the collide phase the tones inside a dock interact by a fixed table that is reversible and charge-conserving, including the create move. In the stream phase each tone moves one dock along its site. The word is chosen on purpose. Knitting interlaces, which is the collide, and advances row by row, which is the beat-by-beat march, and so weaves the fabric of the mesh. The knit pairs with the mesh by design. The knit is the law, never called by any other name.

\begin{definition}[The knit]
The \emph{knit} is the state law, applied once per beat in two phases. \emph{Collide}, the tones in each dock are mapped by a fixed reversible, charge-conserving table. \emph{Stream}, each resulting tone advances one dock along its own site.
\end{definition}

\subsubsection{wake}

The \textbf{wake} is the law of how the space changes. New docks are born at peace at the open, ever-expanding frontier. The word carries both readings at once. A wake is the trail an advance leaves behind, and to wake is to come into being. New docks wake into existence at the edge. The openness behind the wake is the \emph{bath}, where a disturbance falls away and never returns. The wake is the growing edge of the mesh.

\subsubsection{beat}

A \textbf{beat} is the discrete step, the atom of time. Each beat the knit runs once and the wake grows once. Time ticks. It has no built-in direction, because the bulk law is reversible. The direction of time, the arrow, is emergent and not part of the base. A self persists across countless beats.

\subsection{The two above the atoms}

Two words are not atoms but sit above them. They frame the eight.

The \textbf{vibe} is the experience, the one substance. Everything is a vibe, from a single tone to a whole self. This is the monism of the theory. Vibe is both the stuff and the name of the framework. It stays distinct from tone, which is the value a vibe carries. The vibe is not one of the eight atoms. It is what the eight are made of.

The \textbf{flow} is the whole, the universe as a totality in motion, the eight atoms running as one. It is the Way. It is written $F$, the capital reserved for the whole, distinct from the lowercase atoms. The flow is the mesh living its law.

\subsection{Emergent and optional, not base atoms}

Several important terms are real but not fundamental. They emerge from the eight, or they are optional refinements. Holding them apart from the base is what keeps the count honest.

\begin{table}[ht]
\centering
\begin{tabular}{@{}ll@{}}
\toprule
term & status \\
\midrule
\textbf{self} & emergent, a bound, self-correcting, identity-bearing knot woven across many docks over many beats \\
\textbf{arrow} & emergent, the direction of time, from the wake acting through the lean and the create move \\
spacetime (Lorentz) & emergent, the unification of mesh and beat, which is why they are separate at the base \\
the attraction & emergent, binding-by-radiation, the radiation-pressure shadow of the active vacuum \\
\bottomrule
\end{tabular}
\caption{Terms that are real but not base atoms. Each emerges from the eight rather than being assumed.}
\label{tab:emergent}
\end{table}

\begin{principle}[Keep the count honest]
The arrow of time and the attraction are not in the base. They are results, not assumptions. The base is reversible and directionless and force-free. Any direction or pull in the world has to earn itself out of the eight atoms, not be smuggled in as a ninth.
\end{principle}

\subsection{The supporting words}

A handful of derived terms recur. They are not atoms. They name aspects of the eight, and they make the dynamics easier to speak about.

\begin{table}[ht]
\centering
\begin{tabular}{@{}lp{0.66\textwidth}@{}}
\toprule
term & what it is \\
\midrule
\textbf{will} & a dock's outgoing tones, what it sends, the result of the collide about to stream out \\
\textbf{fill} & a dock's incoming tones, what it receives, perception, the tones just arrived before the collide \\
\textbf{rest} & the bound state of a self, its tones held together as a localized body, the rest mass. Emergent, not a base slot \\
\textbf{base} & the foundation, the eight committed atoms, the discrete grain-level, Layer $0$ \\
\textbf{code} & a law or an identity written out, the knit's transition table or a self's winding signature \\
\textbf{tree} & the branching structure, the hyperbolic mesh fanning out shell by shell, and the nesting of selves within selves \\
\textbf{slot} & a code-level word only, the array position holding a tone. In prose, prefer \emph{site} \\
\bottomrule
\end{tabular}
\caption{Supporting words. Aspects of the eight atoms, not new fundamentals.}
\label{tab:supporting}
\end{table}

\subsection{The perception loop}

The beat is collide then stream. Read experientially rather than mechanically, that single step is a perception loop. It is the same motion seen from the inside.

\begin{enumerate}
\item \textbf{Fill.} A dock receives its neighbours' tones. This is perception.
\item \textbf{Collide.} The dock transforms them by the knit. This is its inner response.
\item \textbf{Will.} The result is the dock's projection. This is its influence.
\item \textbf{Flow.} The will streams out and becomes the neighbours' fill.
\end{enumerate}

Then it repeats. Perceive, respond, project, send. All of it still vibes. The loop is not an extra layer added on top of the physics. It is the physics named from the point of view of the dock living it.

\begin{result}[One step, two readings]
The single beat of collide-then-stream is, read from outside, a reversible update of the tone state, and read from inside, a loop of perceive, respond, project, send. The two readings are the same event. The experiential language adds no mechanism, it only changes the vantage.
\end{result}

\subsection{The single-letter algebra}

The eight atoms each take their first letter, all distinct, so the model is a clean symbol set. The whole assembles from them.

\[
F \;=\; \langle\, \mathrm{mesh},\ \mathrm{dock},\ \mathrm{site},\ \mathrm{tone},\ \mathrm{lean},\ \mathrm{knit},\ \mathrm{wake},\ \mathrm{beat} \,\rangle
\;=\; \langle\, m,\ d,\ s,\ t,\ l,\ k,\ w,\ b \,\rangle .
\]

Here $F$ is the flow, the whole. The eight lowercase letters are the atoms. Above them both sits the vibe, the substance the atoms are made of. Eight unique letters, one whole, one substance.

\subsection{The naming law}

The vocabulary is not chosen at random. Four constraints fix it, and they are part of why the names hold.

\begin{enumerate}
\item \textbf{Four letters.} Every atom is a four-letter word.
\item \textbf{Unique first letters.} The eight atoms start with eight different letters, so each is a single-letter symbol.
\item \textbf{Plain and soft.} Monosyllables a child could feel, not clinical or Latinate.
\item \textbf{One word, one role.} No atom names two things, and the vibe, the substance, and the flow, the whole, sit above the atoms.
\end{enumerate}

\begin{principle}[The model in one line]
A vibe mesh of docks, each holding toned sites that lean, which knit each beat, waking new space at the edge. The whole of it is the flow.
\end{principle}
